\begin{abstract}
We prove that stable linearization at one physical H\'enon orbit does not
control the all-conjugate Galois excess assigned to that orbit.  The
distinction matters because a natural continuation of a finite-memory
obstruction ladder is to estimate its discrepancies from the negative fixed
point approached by the physical reflection chains.  We derive instead the
complete reflection-reduced trace algebra at periods eight and nine.  The
period-eight vertex--vertex and edge--edge closures produce irreducible
totally real trace fields of degrees $12$ and $6$; the period-nine
vertex--edge closure produces one irreducible totally real degree-$28$ field
whose extreme real embeddings are the two physical cycles.  Rational root
boxes, modular irreducibility and Sturm counts make these statements exact.
Integer-product bounds then certify two consecutive incidence discrepancies,
\[
\Delta_6=-185.5524168765\ldots<0,
\qquad
\Delta_7=300.0665139420\ldots>0.
\]
The selected physical tail is exponentially localized by a positive stable
multiplier, but the Galois excess sums every nonphysical embedding.  Hence a
physical-tail theorem supplies no bound for the target without a separate,
uniform reflection-ensemble count or pressure theorem.  The result is an
interface obstruction, not an eventual sign law or a Hilbert--P\'olya
construction.
\end{abstract}
